% (c) 2026 Nathanael Fässler, OST Ostschweizer Fachhochschule

\begin{tikzpicture}[scale=1]
  \begin{scope}
  \draw (0,0) grid(4,4);
  \draw[very thick, scale=2] (0,0) grid (2, 2);

  \setcounter{sudokuRow}{1}
    % \setrow {4}{1}  {2}{3}
    % \setrow {2}{3}  {1}{4}

    % \setrow {3}{2}  {4}{1}
    % \setrow {1}{4}  {3}{2}
    % Die "Lösungs-Zahlen" sind blau gefärbt und fett.
    \setrow {\textcolor{blue!80!black}{\textbf{4}}}{\textcolor{blue!80!black}{\textbf{1}}} {2}{\textcolor{blue!80!black}{\textbf{3}}}
    \setrow {\textcolor{blue!80!black}{\textbf{2}}}{\textcolor{blue!80!black}{\textbf{3}}} {\textcolor{blue!80!black}{\textbf{1}}}{\textcolor{blue!80!black}{\textbf{4}}}

    \setrow {3}{\textcolor{blue!80!black}{\textbf{2}}} {\textcolor{blue!80!black}{\textbf{4}}}{1}
    \setrow {\textcolor{blue!80!black}{\textbf{1}}}{4} {\textcolor{blue!80!black}{\textbf{3}}}{\textcolor{blue!80!black}{\textbf{2}}}
    
    \node[anchor=center] at (2, -0.5) {};
  \end{scope}

\end{tikzpicture}
